%%%%%%%%%%%%%%%%%%%%%%%%%%%%%%%%%%%%%%%%%%%%%%%%%%%%%%%%%%%%%%%%%%%%%%%%%%%%%%%%
% Obxectivo: Lista de termos empregados no documento,                          %
%            xunto cos seus respectivos significados.                          %
%%%%%%%%%%%%%%%%%%%%%%%%%%%%%%%%%%%%%%%%%%%%%%%%%%%%%%%%%%%%%%%%%%%%%%%%%%%%%%%%

\newglossaryentry{cnn}{
  name={CNN},
  description={Rede Neuronal Convolucional. Arquitectura de aprendizaxe profunda especialmente eficaz para o procesado de imaxes e sinais mediante capas de convolución.}
}

\newglossaryentry{coral}{
  name={CORAL},
  description={No marco de regresión ordinal para redes neuronais, CORAL é un método que garante monotonía no ranking, puntuación de confianza coherente e independencia da arquitectura.
  }
}

\newglossaryentry{hpi}{
  name={HPI},
  description={Índice de presión de depredación herbívora. É unha medida que avalía o dano por herbívoros que ten unha plante, no noso caso unha alga Laminaria. Un valor baixo (próximo a 0) indica pouco dano, e un valor alto (próximo a 6) indica danos maiores á alga)}
}

\newglossaryentry{ivr}{
  name={IVR},
  description={Ratio de vertebrado vs invertebrado. Cuantifica, entre vertebrados e invertebrados, a seguridade de quén causou os danos. \textcolor{red}{\textbf{TODO: Qué es valor bajo y qué es valor alto??.}}}
}

\newglossaryentry{kelp}{
  name={Kelp / \textit{Laminaria}},
  description={\textit{Kelp}: nome común de grandes algas laminarias en inglés. \textit{Laminaria} é un xénero típico de algas en mares fríos e temperados.}
}

\newglossaryentry{yolo}{
  name={YOLO},
  description={É unha serie de sistemas de detección de obxectos en tempo real baseado en \glspl{cnn}. Esta tecnoloxía recibiu melloras ata converterse no \textit{framework} máis popular de detección de obxectos.}
}
%. is a series of real-time object detection systems based on convolutional neural networks
 % YOLO has undergone several iterations and improvements, becoming one of the most popular object detection frameworks.

%TODO: BORRAR ESTE DUMMY
\newglossaryentry{bytecode}{
  name=bytecode,
  description={Código independente da máquina que xeran compiladores de determinadas linguaxes (Java, Erlang,\dots) e que é executado polo correspondente intérprete.}
}